% SEGMENT OLP-0025-S01
% Part: sets-functions-relations
% Chapter: functions
% Section: composition

\documentclass[../../../include/open-logic-section]{subfiles}

\begin{document}

\olfileid{sfr}{fun}{cmp}
\olsection{சார்புகளின் சேர்ப்பு}

\begin{explain}
\oliflabeldef{sfr:fun:inv:sec}{!!a{bijection} $f$ இன்
நேர்மாறு $f^{-1}$ உம் ஒரு சார்பே என்பதை \olref[inv]{sec} இல் கண்டோம்.
சார்புகளின் மீதான மற்றொரு செயல் சேர்ப்பு ஆகும்.}{}
$f$, $g$ ஆகிய இரண்டு சார்புகளைச் சேர்ப்பதன் மூலம்,
அதாவது முதலில் $f$ ஐயும் பின்னர் $g$ ஐயும் செயல்படுத்துவதன் மூலம்,
ஒரு புதிய சார்பை வரையறுக்கலாம்.
வீச்சகங்களும் சார்பகங்களும் பொருந்தினால் மட்டுமே இது சாத்தியம்:
அதாவது, $f$ இன் வீச்சகம் $g$ இன் சார்பகத்தின் ஓர் உட்கணமாக
இருக்க வேண்டும்.
\oliflabeldef{sfr:rel:ops:sec}{சார்புகளின் மீதான இந்தச் செயல்,
\olref[rel][ops]{sec} இல் கண்ட தொடர்புகளின் சார்புப் பெருக்கம்
என்ற செயலுக்கு ஒப்பானதாகும்.}{}

சேர்ப்பு என்ற கருத்தை விளக்க ஒரு படம் உதவும்.
\olref{fig:composition} இல் $f \colon A \to B$,
$g \colon B \to C$ ஆகிய இரண்டு சார்புகளையும் அவற்றின்
சேர்ப்பான $(\comp{f}{g})$ ஐயும் காட்டுகிறோம்.
$(\comp{f}{g}) \colon A \to C$ என்ற சார்பு,
$A$ இன் ஒவ்வோர் !!{element}[acc] $C$ இன் !!a{element} உடன்
இணைக்கிறது. $A$ இன் !!a{element} ஆனது $C$ இன் எந்த
!!{element} உடன் இணைக்கப்படுகிறது என்பதைப் பின்வருமாறு
குறிப்பிடுகிறோம்: $x \in
A$ என்ற உள்ளீடு தரப்பட்டால், முதலில் $x$ மீது $f$ என்ற
சார்பைச் செயல்படுத்துகிறோம். அது ஏதாவது ஒரு
$f(x)
= y \in B$ ஐ வெளியீடாகத் தருகிறது.
பின்னர் $y$ மீது $g$ என்ற சார்பைச் செயல்படுத்துகிறோம்.
அது ஏதாவது ஒரு $g(f(x)) = g(y) = z \in C$ ஐ வெளியீடாகத் தருகிறது.
\begin{figure}
  \olasset[2\olphotowidth]{assets/diagrams/composition.tikz}
  \caption{$f$, $g$ ஆகிய இரண்டு சார்புகளின் சேர்ப்பு $g \circ f$.}
  \ollabel{fig:composition}
\end{figure}
\end{explain}

% SEGMENT OLP-0025-S02
\begin{defn}[சேர்ப்பு]
$f\colon A \to B$, $g\colon B \to C$ ஆகியவை சார்புகளாக இருக்கட்டும்.
$f$ ஐ $g$ உடன் \emph{சேர்ப்பதால்} கிடைப்பது
$\comp{f}{g} \colon A \to C$ ஆகும்;
இங்கு $(\comp{f}{g})(x) = g(f(x))$.
\end{defn}

% SEGMENT OLP-0025-S03
\begin{ex}
$f(x) = x + 1$, $g(x) = 2x$ ஆகிய சார்புகளைக் கருதுவோம்.
$(\comp{f}{g})(x) = g(f(x))$ என்பதால்,
ஒவ்வோர் உள்ளீடு $x$ க்கும் முதலில் அதன் தொடரியை எடுத்துக்கொண்டு,
பின்னர் கிடைப்பதை இரண்டால் பெருக்க வேண்டும்.
ஆகவே அவற்றின் சேர்ப்பு $(\comp{f}{g})(x) = 2(x+1)$ எனக்
கொடுக்கப்படுகிறது.
\end{ex}

% SEGMENT OLP-0025-S04
\begin{prob}
$f \colon A \to B$, $g \colon B \to C$ ஆகிய இரண்டும்
!!{injective} பண்புடையவை எனில்,
$\comp{f}{g}\colon A \to C$ என்பதும் !!{injective} பண்புடையது
எனக் காட்டுக.
\end{prob}

% SEGMENT OLP-0025-S05
\begin{prob}
$f \colon A \to B$, $g \colon B \to C$ ஆகிய இரண்டும்
!!{surjective} பண்புடையவை எனில்,
$\comp{f}{g}\colon A \to C$ என்பதும் !!{surjective} பண்புடையது
எனக் காட்டுக.
\end{prob}

% SEGMENT OLP-0025-S06
\begin{prob}
$f \colon A \to B$, $g \colon B \to C$ எனக் கொள்வோம்.
$\comp{f}{g}$ இன் கோட்டுரு $R_f \mid R_g$ எனக் காட்டுக.
\end{prob}

\end{document}
